% book_geometry/sections/section_01_angles.tex
\chapter{Кути. Суміжні та вертикальні кути}

\section{Попередні поняття}

\paragraph{13. Кут.}
Фігура, утворена двома променями ($OA$ і $OB$, рис.~8), що виходять з однієї точки, називається \textbf{кутом}. Промені, які утворюють кут, називаються \textbf{сторонами кута}, а точка, з якої вони виходять, — \textbf{вершиною кута}. Сторони кута слід уявляти собі нескінченно продовженими від вершини.

\begin{center}
\begin{tikzpicture}[scale=1.1, >=Stealth]
    \coordinate (O) at (0,0);
    \coordinate (A) at (4,2.5);
    \coordinate (B) at (4.5,0);
    \coordinate (D) at (3.8,1.6);
    \coordinate (E) at (4.2,0.8);

    \draw[->, very thick, color=primaryblue] (O) -- (A) node[above right] {$A$};
    \draw[->, very thick, color=primaryblue] (O) -- (B) node[right] {$B$};
    
    \draw[->, thick, dashed, color=gray] (O) -- (D) node[right] {$D$};
    \draw[->, thick, dashed, color=gray] (O) -- (E) node[right] {$E$};

    \filldraw[black] (O) circle (2pt) node[left] {$O$};
    \node at (2.2, -0.6) {\small\textbf{Рис. 8.} Кут $AOB$ та внутрішні промені $OD, OE$};
\end{tikzpicture}
\end{center}

Кут зазвичай позначають трьома буквами, з яких середня ставиться біля вершини, а крайні — біля будь-яких точок на сторонах; наприклад, кажуть: «кут $AOB$» або «кут $BOA$» (рис.~8). Проте можна позначати кут і однією буквою, поставленою біля вершини, якщо при цій вершині немає інших кутів. Іноді ми позначатимемо кут цифрою, поставленою всередині кута біля вершини.

Частина площини, обмежена сторонами кута, розглядається як така, що лежить \textbf{всередині кута}; решта площини лежить \textbf{поза кутом}; перша становить \textbf{внутрішню область кута}, друга — \textbf{зовнішню область} його.

Якщо з вершини кута (рис.~8) проведемо всередині нього будь-які промені $OD, OE \dots$, то утворені при цьому кути $AOD, DOE, EOB \dots$ розглядаються як частини кута $AOB$.

Слово «кут» на письмі часто замінюють символом $\angle$.

\begin{feynmanbox}[Кут — це міра повороту, а не просто лінії на папері]
Коли ми дивимося на креслення в зошиті, кут здається застиглим трикутним шматочком площини. Проте в реальному фізичному світі кут — це \textbf{динамічний процес обертання}.

Уяви стрілку годинника, циркуль чи промінь обертового маяка. Спочатку промінь був спрямований уздовж рейки $OB$. Потім його повернули навколо шарніра $O$ до положення $OA$. 

Величина кута показує, \textbf{наскільки сильно повернувся промінь}. Самі промені нескінченні, тому їхня довжина на папері ніяк не впливає на кут!
\end{feynmanbox}

\motionlink{Кут як поворот}%
{Повертай промінь навколо вершини й спостерігай, як положення променя визначає величину кута.}%
{https://radchr.github.io/Algebra/geometry-angle-rotation.html}%
{qr_angle_rotation.png}

\paragraph{14. Рівність і нерівність кутів.}
Два кути вважаються \textbf{рівними}, якщо при накладанні вони можуть суміститися. Припустимо, наприклад, що ми накладаємо кут $AOB$ на кут $A_1O_1B_1$ (рис.~9) так, щоб вершина $O$ збіглася з $O_1$, сторона $OB$ пішла по $O_1B_1$ і щоб кути накрили один одного своїми внутрішніми областями. Якщо при цьому сторона $OA$ суміститься з $O_1A_1$, то кути рівні; якщо ж сторона $OA$ пройде всередині кута $A_1O_1B_1$ або поза ним, то кути не рівні, до того ж меншим буде той з них, який становить частину іншого кута.

\begin{center}
\begin{tikzpicture}[scale=0.95, >=Stealth]
    % Перший кут AOB
    \coordinate (O) at (0,0);
    \coordinate (A) at (2.4,2.1);
    \coordinate (B) at (3.2,0);
    \draw[->, very thick, color=primaryblue] (O) -- (A) node[above left] {$A$};
    \draw[->, very thick, color=primaryblue] (O) -- (B) node[right] {$B$};
    \filldraw (O) circle (2pt) node[left] {$O$};

    % Другий кут A1O1B1 з накладанням
    \coordinate (O1) at (5.5,0);
    \coordinate (A1) at (7.9,2.1);
    \coordinate (B1) at (8.7,0);
    \coordinate (A_in) at (7.6,1.4);
    \draw[->, very thick, color=primaryblue] (O1) -- (A1) node[above right] {$A_1$};
    \draw[->, very thick, color=primaryblue] (O1) -- (B1) node[right] {$B_1$};
    \draw[->, thick, dashed, color=warningred] (O1) -- (A_in) node[right] {\small $OA$ (всередині)};
    \filldraw (O1) circle (2pt) node[below left] {$O_1$};

    \node at (4.2, -0.8) {\small\textbf{Рис. 9.} Порівняння кутів способом накладання};
\end{tikzpicture}
\end{center}

\motionlink{Рівні кути можна сумістити}%
{Перенеси й поверни копію кута до точного суміщення з еталоном та простеж, що його величина не змінюється.}%
{https://radchr.github.io/Algebra/geometry-equal-angles-motion.html}%
{qr_equal_angles_motion.png}

\paragraph{15. Сума кутів.}
\textbf{Сумою кутів} $AOB$ і $A_1O_1B_1$ (рис.~10) називається такий кут, який отримують у такий спосіб. Будуємо кут $MNP$, що дорівнює першому даному куту $AOB$, і до нього приєднуємо кут $PNQ$, що дорівнює другому даному куту $A_1O_1B_1$, так, щоб в обох кутів виявилася спільна вершина $N$ і спільна сторона $NP$, а внутрішні області кутів були розташовані по різні боки від цієї спільної сторони $NP$. Тоді кут $MNQ$ і буде сумою кутів $AOB$ та $A_1O_1B_1$. Подібним чином може бути складена сума трьох і більше кутів.

\begin{center}
\begin{tikzpicture}[scale=0.95, >=Stealth]
    % Кут 1
    \draw[->, thick] (0,0) node[left] {$O$} -- (2,1.5) node[above] {$B$};
    \draw[->, thick] (0,0) -- (2.3,0) node[right] {$A$};
    \draw (0.6,0) arc (0:37:0.6);

    % Кут 2
    \draw[->, thick] (3.5,0) node[left] {$O_1$} -- (5.5,1.0) node[above] {$B_1$};
    \draw[->, thick] (3.5,0) -- (5.8,0) node[right] {$A_1$};
    \draw (4.1,0) arc (0:26:0.6);

    % Сума
    \coordinate (N) at (7.5,0);
    \draw[->, very thick, color=primaryblue] (N) node[below left] {$N$} -- (10.2,0) node[right] {$M$};
    \draw[->, thick, dashed] (N) -- (9.8,1.6) node[above right] {$P$};
    \draw[->, very thick, color=primaryblue] (N) -- (9.0,2.4) node[above] {$Q$};
    \draw[color=primaryblue] (8.2,0) arc (0:58:0.7);

    \node at (5.5, -0.8) {\small\textbf{Рис. 10.} Додавання двох кутів: $\angle MNQ = \angle AOB + \angle A_1O_1B_1$};
\end{tikzpicture}
\end{center}

Сума кутів, як і сума відрізків прямої, має \textbf{переставну} та \textbf{сполучну} властивості.

З поняття про суму кутів виводяться поняття про їхню \textbf{різницю}, \textbf{добуток} і \textbf{частку}.

Часто доводиться говорити про такий промінь, який ділить даний кут навпіл; такому променю дали спеціальну назву: \textbf{бісектриса} (рис.~11).

\begin{center}
\begin{tikzpicture}[scale=1.0, >=Stealth]
    \coordinate (O) at (0,0);
    \coordinate (A) at (4,2.4);
    \coordinate (B) at (4.5,0);
    \coordinate (L) at (4.4,1.2);

    \draw[->, thick] (O) -- (A) node[above right] {$A$};
    \draw[->, thick] (O) -- (B) node[right] {$B$};
    \draw[->, very thick, color=feynmangreen] (O) -- (L) node[right] {\textbf{Бісектриса $OL$}};

    \draw[color=feynmangreen] (0.9,0) arc (0:15.5:0.9);
    \draw[color=feynmangreen] (0.9,0.25) arc (15.5:31:0.9);
    \filldraw[black] (O) circle (2pt) node[left] {$O$};
    \node at (2.2, -0.6) {\small\textbf{Рис. 11.} Бісектриса кута};
\end{tikzpicture}
\end{center}

\paragraph{16. Розширення поняття про кут.}
Під час знаходження суми кутів можуть трапитися деякі особливі випадки, які корисно розглянути окремо:

\begin{enumerate}
    \item Може статися, що після додавання кількох кутів, наприклад трьох: $AOB, BOC$ і $COD$ (рис.~12), сторона $OD$ кута $COD$ стане продовженням сторони $OA$ кута $AOB$. Ми отримаємо тоді фігуру, утворену двома променями ($OA$ і $OD$), що виходять з однієї точки ($O$) і становлять продовження один одного (є доповняльними променями). Таку фігуру прийнято теж називати кутом (\textbf{розгорнутим}). Внутрішня область розгорнутого кута становить половину площини.
    
    \begin{center}
    \begin{tikzpicture}[scale=1.0, >=Stealth]
        \coordinate (O) at (0,0);
        \coordinate (A) at (3,0);
        \coordinate (B) at (1.5,1.8);
        \coordinate (C) at (-1.2,1.8);
        \coordinate (D) at (-3,0);

        \draw[->, thick] (O) -- (A) node[right] {$A$};
        \draw[->, thick] (O) -- (B) node[above right] {$B$};
        \draw[->, thick] (O) -- (C) node[above left] {$C$};
        \draw[->, thick] (O) -- (D) node[left] {$D$};
        
        \draw[color=primaryblue, thick] (0.8,0) arc (0:180:0.8);
        \filldraw (O) circle (2pt) node[below] {$O$};
        \node at (0, -0.6) {\small\textbf{Рис. 12.} Розгорнутий кут як сума суміжних частин};
    \end{tikzpicture}
    \end{center}

    \item Може статися, що після додавання кількох кутів, наприклад п'яти кутів: $AOB,\allowbreak BOC,\allowbreak COD,\allowbreak DOE$ і $EOA$ (рис.~13), сторона $OA$ кута $EOA$ суміститься зі стороною $OA$ кута $AOB$. Фігура, утворена такими суміщеними променями (розглядувана разом з усією площиною, розташованою навколо їхньої спільної вершини $O$), також називається кутом (\textbf{повним кутом}).
    
    \begin{center}
    \begin{tikzpicture}[scale=0.9, >=Stealth]
        \coordinate (O) at (0,0);
        \draw[->, thick] (O) -- (2.5,0) node[right] {$A$};
        \draw[->, thick] (O) -- (1.5,2.0) node[above right] {$B$};
        \draw[->, thick] (O) -- (-1.2,2.0) node[above left] {$C$};
        \draw[->, thick] (O) -- (-2.0,-0.8) node[left] {$D$};
        \draw[->, thick] (O) -- (0.5,-2.0) node[below] {$E$};

        \draw[color=primaryblue, thick, ->] (0.6,0) arc (0:350:0.6);
        \filldraw (O) circle (2pt) node[below left=2pt] {$O$};
        \node at (0, -2.6) {\small\textbf{Рис. 13.} Повний кут навколо точки $O$};
    \end{tikzpicture}
    \end{center}

    \item Нарешті, може статися, що, будуючи суму кутів, ми не лише заповнимо всю площину навколо їхньої спільної вершини, а й навіть будемо змушені накладати кути один на другий, покриваючи площину навколо спільної вершини вдруге, втретє тощо. Така сума кутів дорівнює одному повному куту, складеному з деяким кутом, або дорівнює двом повним кутам, складеним з деяким кутом, тощо.
\end{enumerate}

\begin{warningbox}[Розгорнутий кут — це не просто пряма лінія!]
Часто семикласники вважають: «Розгорнутий кут — це звичайна пряма».

Це не так: у прямої лінії немає виділеного центра. А розгорнутий кут має чітко зафіксовану \textbf{вершину $O$} та два протилежні промені. Кут розглядається разом із внутрішньою областю (півплощиною) з одного боку від прямої!
\end{warningbox}

---

\section{Вимірювання кутів}

\paragraph{17. Центральний кут і два його властивості.}
Кут ($AOB$, рис.~14), утворений двома радіусами, називається \textbf{центральним кутом}; про такий кут і дугу, укладену між його сторонами, кажуть, що вони відповідають один одному.

\begin{center}
\begin{tikzpicture}[scale=0.9, >=Stealth]
    % Рис. 14
    \draw[thick] (0,0) circle (1.5);
    \coordinate (O) at (0,0);
    \coordinate (A) at (-1.1,-1.0);
    \coordinate (B) at (1.1,-1.0);
    \draw[thick] (O) -- (A) node[below left] {$A$};
    \draw[thick] (O) -- (B) node[below right] {$B$};
    \filldraw (O) circle (1.5pt) node[above] {$O$};
    \node at (0, -1.9) {\small\textbf{Рис. 14.} Центральний кут};

    % Рис. 15
    \begin{scope}[xshift=6.5cm]
        \draw[thick] (0,0) circle (1.5);
        \coordinate (O2) at (0,0);
        \coordinate (A2) at (-0.6,1.4);
        \coordinate (B2) at (1.2,0.9);
        \coordinate (C2) at (1.4,-0.4);
        \coordinate (D2) at (0.2,-1.5);
        \draw[thick] (O2) -- (A2) node[above left] {$A$};
        \draw[thick] (O2) -- (B2) node[above right] {$B$};
        \draw[thick] (O2) -- (C2) node[right] {$C$};
        \draw[thick] (O2) -- (D2) node[below right] {$D$};
        \draw[->, thick, color=mediapurple] (1.7,0.7) arc (20:-15:1.7);
        \filldraw (O2) circle (1.5pt) node[left] {$O$};
        \node at (0, -1.9) {\small\textbf{Рис. 15.} Рівні центральні кути};
    \end{scope}
\end{tikzpicture}
\end{center}

Центральні кути відносно відповідних їм дуг мають такі дві властивості:

В одному колі або в рівних колах:
\begin{enumerate}
    \item \textbf{Якщо центральні кути рівні, то й відповідні їм дуги рівні.}
    \item \textbf{Обернено: якщо дуги рівні, то й відповідні їм центральні кути рівні.}
\end{enumerate}

\begin{proof}
Нехай $\angle AOB = \angle COD$ (рис.~15); потрібно довести, що дуги $AB$ і $CD$ також рівні.

Уявімо, що сектор $AOB$ ми повернули навколо центра $O$ у напрямку, вказаному стрілкою, настільки, щоб радіус $OA$ збігся з $OC$. Тоді, внаслідок рівності кутів, радіус $OB$ суміститься з $OD$; значить, сумістяться й дуги $AB$ та $CD$, тобто вони будуть рівні.

Друга властивість доводиться так само накладанням.
\end{proof}

\motionlink{Центральний кут і відповідна дуга}%
{Змінюй два рівні центральні кути в рівних колах і спостерігай, чому відповідні їм дуги також залишаються рівними.}%
{https://radchr.github.io/Algebra/geometry-central-angle-arc.html}%
{qr_central_angle_arc.png}

\paragraph{18. Градуси: дуговий і кутовий.}
Уявімо, що деяке коло поділено на 360 рівних частин і з кожної точки поділу проведено радіуси до центра. Тоді навколо центра утворяться 360 маленьких центральних кутів, які, як відповідні рівним дугам згідно з теоремою попереднього параграфа, мають бути рівними між собою.

Кожна з отриманих у такий спосіб на колі маленьких дуг називається \textbf{дуговим градусом}, а кожен з утворених при центрі маленьких кутів називається \textbf{кутовим градусом}.

Отже, можна сказати, що:
\begin{itemize}
    \item \textbf{дуговий градус} є $\frac{1}{360}$ частиною кола (або $\frac{1}{180}$ частиною півкола, або $\frac{1}{90}$ частиною чверті кола);
    \item \textbf{кутовий градус} є центральним кутом, що відповідає дуговому градусу.
\end{itemize}

Оскільки 360 таких центральних кутів становлять повний кут, то можна також сказати, що кутовий градус є $\frac{1}{360}$ частиною повного кута (або $\frac{1}{180}$ частиною розгорнутого кута, або $\frac{1}{90}$ частиною прямого кута).

Градуси (дугові та кутові) поділяються ще на 60 рівних частин, які називають \textbf{мінутами}, а мінути поділяються ще на 60 рівних частин, які називають \textbf{секундами}.

\paragraph{19. Відповідність між центральними кутами й дугами.}
Нехай $AOB$ є довільний кут (рис.~16). Опишемо між його сторонами з вершини $O$ як із центра довільним радіусом дугу $CD$; тоді кут $AOB$ буде центральним кутом, відповідним дузі $CD$.

\begin{center}
\begin{tikzpicture}[scale=1.1, >=Stealth]
    \coordinate (O) at (0,0);
    \draw[thick] (O) -- (4,2.2) node[right] {$A$};
    \draw[thick] (O) -- (4.2,0) node[right] {$B$};
    \draw[very thick, color=primaryblue] (2.8,0) arc (0:29:2.8) node[midway, right=2pt] {\small $CD$};
    \node at (2.9, 1.6) {$C$};
    \node at (2.9, -0.3) {$D$};

    % Радіальні поділки
    \foreach \a in {3, 6, 9, 12, 15, 18, 21, 24, 27} {
        \draw[thin, color=gray] (O) -- (\a:2.8);
    }
    \filldraw (O) circle (2pt) node[left] {$O$};
    \node at (2.2, -0.7) {\small\textbf{Рис. 16.} Вимірювання кута дугою кола};
\end{tikzpicture}
\end{center}

Припустимо, що в цій дузі міститься 20 дугових градусів (на рисунку градуси зображено у збільшеному розмірі). Тоді, якщо сполучимо точки поділу з центром, кут $AOB$ розділиться, очевидно, на 20 кутових градусів.

Узагалі можна сказати, що \textbf{кут вимірюється відповідною йому дугою}, розуміючи під цим те, що в куті міститься стільки кутових градусів, мінут і секунд, скільки у відповідній йому дузі міститься дугових градусів, мінут і секунд. Якщо, наприклад, у дузі $CD$ міститься 20 градусів 10 мінут 15 секунд дугових, то й у куті $AOB$ міститься 20 градусів 10 мінут 15 секунд кутових, що прийнято скорочено записувати так:
\[
\angle AOB = 20^\circ 10' 15'',
\]
позначаючи знаками $(^\circ), ('), ('')$ відповідно градуси, мінути й секунди.

\paragraph{20. Транспортир.}
Цей прилад (рис.~17), що застосовується для вимірювання кутів, являє собою півкруг, дуга якого розділена на 180 градусів. Щоб виміряти кут $DCE$, накладають на нього прилад так, щоб центр півкруга збігся з вершиною кута, а базовий радіус $CB$ збігся зі стороною $CE$. Тоді число градусів, що міститься в дузі, укладеній між сторонами кута $DCE$, покаже його величину.

За допомогою транспортира можна також накреслити кут, що містить задане число градусів (наприклад, кут у $90^\circ, 45^\circ, 30^\circ$ тощо).

\begin{center}
\begin{tikzpicture}[scale=1.0, >=Stealth]
    % Корпус транспортира
    \draw[thick, fill=primaryblue!6] (-3.2,0) -- (3.2,0) arc (0:180:3.2) -- cycle;
    \draw[thick, fill=white] (-2.2,0) -- (2.2,0) arc (0:180:2.2) -- cycle;
    \filldraw (0,0) circle (2pt) node[below] {$C$};
    
    % Градусні поділки
    \foreach \a in {0,10,...,180} {
        \draw (\a:2.8) -- (\a:3.2);
    }
    \foreach \a in {0,30,...,180} {
        \draw[thick] (\a:2.6) -- (\a:3.2);
        \node[font=\tiny] at (\a:2.4) {\a$^\circ$};
    }
    
    % Вимірюваний кут DCE = 50°
    \draw[very thick, color=primaryblue, ->] (0,0) -- (3.5,0) node[right] {$E$};
    \draw[very thick, color=warningred, ->] (0,0) -- (50:3.5) node[above right] {$D$};
    \draw[color=warningred, thick] (1.1,0) arc (0:50:1.1);
    \node[color=warningred, font=\small\bfseries] at (25:1.45) {$50^\circ$};
    
    \node at (0, -0.6) {\small\textbf{Рис. 17.} Вимірювання кута $DCE = 50^\circ$ за допомогою транспортира};
\end{tikzpicture}
\end{center}

\newpage

\section{Суміжні та вертикальні кути}

\paragraph{21. Прямий, гострий і тупий кути.}
Кут у $90^\circ$ (який становить, отже, половину розгорнутого кута або четвертину повного кута) називають \textbf{прямим кутом}; кут, менший від прямого, називають \textbf{гострим}, а кут, більший за прямий, називають \textbf{тупим} (рис.~18).

\begin{center}
\begin{tikzpicture}[scale=0.9, >=Stealth]
    % Прямий кут
    \begin{scope}[xshift=0cm]
        \draw[thick] (0,0) -- (2.2,0);
        \draw[thick] (0,0) -- (0,2.2);
        \draw (0.4,0) -- (0.4,0.4) -- (0,0.4);
        \draw[dashed] (1.5,0) arc (0:90:1.5);
        \node at (0.9,0.9) {$90^\circ$};
        \node at (1.0, -0.5) {\small\textbf{прямий}};
    \end{scope}

    % Гострий кут
    \begin{scope}[xshift=4.5cm]
        \draw[thick] (0,0) -- (2.2,0);
        \draw[thick] (0,0) -- (1.5,2.0);
        \draw[color=primaryblue, thick] (0.8,0) arc (0:53:0.8);
        \node at (1.0, -0.5) {\small\textbf{гострий}};
    \end{scope}

    % Тупий кут
    \begin{scope}[xshift=9.0cm]
        \draw[thick] (0,0) -- (2.2,0);
        \draw[thick] (0,0) -- (-1.2,1.8);
        \draw[color=warningred, thick] (0.8,0) arc (0:124:0.8);
        \node at (0.5, -0.5) {\small\textbf{тупий}};
    \end{scope}

    \node at (5.0, -1.2) {\small\textbf{Рис. 18.} Види кутів за градусною мірою};
\end{tikzpicture}
\end{center}

Звісно, всі прямі кути, як такі, що містять однакове число градусів, рівні між собою. Величину прямого кута в класичних підручниках іноді позначають буквою $d$ (початкова буква французького слова \textit{droit}, що означає «прямий»).

\begin{feynmanbox}[Погляд у майбутнє: Як змінювалася мова кутів — від Евкліда до числа $\pi$]
У математиці один і той самий розгорнутий кут описують по-різному залежно від епохи та розділу науки:
\begin{itemize}
    \item \textbf{В античності (Евклід, IV ст. до н.~е.):} градусів у геометрії взагалі не було. Головним природним еталоном вважався прямий кут. Тому розгорнутий кут називали «сумою двох прямих кутів» (скорочено $2d$, від фр.~\textit{droit} — прямий, як писав Кисельов).
    \item \textbf{У 7 класі (градусна міра):} ми користуємося градусами ($180^\circ$) — це практичний спадок стародавніх вавилонських астрономів, які пов'язували поділ кола на 360 частин із рухом Сонця небосхилом протягом року.
    \item \textbf{У старших класах та фізиці (радіани та число $\pi$):} у 9--10 класах ви дізнаєтеся про фундаментальну міру — \textbf{радіан}. Виявиться, що розгорнутий кут дорівнює точно \textbf{$\pi$ радіан} ($\approx 3{,}14159\dots$). Чому саме $\pi$? Бо якщо накреслити коло радіусом 1 метр, то довжина дуги розгорнутого кута (рівно половина цього кола!) становитиме саме $\pi$ метрів!
\end{itemize}

Запам'ятай цей красивий ланцюжок еволюції геометричної мови:
\[
\text{\textbf{Розгорнутий кут}} = \underbrace{\mathbf{2d}}_{\text{Евклід (еталон)}} = \underbrace{\mathbf{180^\circ}}_{\text{7 клас (градуси)}} = \underbrace{\boldsymbol{\pi}\text{ \textbf{рад}}}_{\text{старша школа (число } \pi)}.
\]
\end{feynmanbox}

\paragraph{22. Суміжні кути та їхні властивості.}
Два кути ($AOB$ і $BOC$, рис.~19) називаються \textbf{суміжними}, якщо одна сторона в них спільна, а дві інші сторони становлять продовження одна одної (є доповняльними променями).

\begin{center}
\begin{tikzpicture}[scale=1.1, >=Stealth]
    \coordinate (O) at (0,0);
    \coordinate (A) at (-3,0);
    \coordinate (C) at (3,0);
    \coordinate (B) at (1.5,2.2);

    \draw[thick] (A) -- (C);
    \draw[->, very thick, color=primaryblue] (O) -- (B) node[above right] {$B$};
    
    \draw[color=primaryblue, thick] (0.7,0) arc (0:56:0.7) node[midway, right=2pt] {$\beta$};
    \draw[color=warningred, thick] (-0.7,0) arc (180:56:0.7) node[midway, left=2pt] {$\alpha$};

    \filldraw (A) circle (1.5pt) node[below] {$A$};
    \filldraw (O) circle (2pt) node[below] {$O$};
    \filldraw (C) circle (1.5pt) node[below] {$C$};

    \node at (0, -0.7) {\small\textbf{Рис. 19.} Суміжні кути $\angle AOB$ ($\alpha$) та $\angle BOC$ ($\beta$)};
\end{tikzpicture}
\end{center}

Оскільки такі кути в сумі становлять розгорнутий кут, то \textbf{сума двох суміжних кутів дорівнює $180^\circ$} (іншими словами, вона дорівнює сумі двох прямих кутів: $2d$).

Для кожного даного кута можна побудувати два суміжні з ним кути. Наприклад, для кута $AOB$ (рис.~20), продовживши сторону $AO$, ми отримаємо один суміжний кут $BOC$, а продовживши сторону $BO$, отримаємо другий суміжний кут $AOD$. 

\begin{center}
\begin{tikzpicture}[scale=1.0, >=Stealth]
    \coordinate (O) at (0,0);
    \coordinate (A) at (3,0);
    \coordinate (C) at (-2.5,0);
    \coordinate (B) at (1.8,2.2);
    \coordinate (D) at (-1.8,-2.2);

    \draw[thick] (C) node[left] {$C$} -- (A) node[right] {$A$};
    \draw[thick] (D) node[left] {$D$} -- (B) node[right] {$B$};
    \filldraw (O) circle (2pt) node[below right=2pt] {$O$};

    \node at (1.2, 0.4) {$\angle AOB$};
    \node at (-0.8, 0.5) {$\angle BOC$};
    \node at (0.3, -0.8) {$\angle AOD$};

    \node at (0, -2.8) {\small\textbf{Рис. 20.} Два кути ($BOC$ і $AOD$), суміжні з кутом $AOB$};
\end{tikzpicture}
\end{center}

Два кути, суміжні з одним і тим самим кутом, рівні між собою, тому що кожен із них містить однакове число градусів, а саме таке число, яке в сумі з числом градусів кута $AOB$ становить $180^\circ$.

Якщо кут $AOB$ прямий (рис.~21), тобто містить $90^\circ$, то й кожен із суміжних із ним кутів $COB$ та $AOD$ має бути також прямим, оскільки містить $180^\circ - 90^\circ = 90^\circ$; четвертий кут $COD$ теж має бути прямим, оскільки три кути $AOB, BOC$ і $AOD$ становлять у сумі $270^\circ$, і отже, від $360^\circ$ на долю кута $COD$ залишається теж $90^\circ$.

Таким чином: \textbf{якщо при перетині двох прямих ($AC$ і $BD$, рис.~21) один із чотирьох кутів виявиться прямим, то й решта три кути мають бути прямими.}

\begin{center}
\begin{tikzpicture}[scale=0.9, >=Stealth]
    \draw[thick, <->] (-2.5,0) -- (2.5,0) node[right] {$A$};
    \draw[thick, <->] (0,-2.5) -- (0,2.5) node[above] {$B$};
    \node at (-2.7, 0) {$C$};
    \node at (0, -2.7) {$D$};
    \filldraw (0,0) circle (2pt) node[below left] {$O$};

    \draw (0.4,0) -- (0.4,0.4) -- (0,0.4);
    \draw (-0.4,0) -- (-0.4,0.4) -- (0,0.4);
    \draw (0.4,0) -- (0.4,-0.4) -- (0,-0.4);
    \draw (-0.4,0) -- (-0.4,-0.4) -- (0,-0.4);

    \node at (0.8,0.8) {$90^\circ$};
    \node at (-0.8,0.8) {$90^\circ$};
    \node at (-0.8,-0.8) {$90^\circ$};
    \node at (0.8,-0.8) {$90^\circ$};
    \node at (0, -3.2) {\small\textbf{Рис. 21.} Перетин прямих під прямим кутом: $AC \perp BD$};
\end{tikzpicture}
\end{center}

\motionlink{Суміжні кути}%
{Рухай спільний промінь і перевіряй: один кут зростає рівно настільки, наскільки другий зменшується, тому їхня сума залишається $180^\circ$.}%
{https://radchr.github.io/Algebra/geometry-adjacent-angles.html}%
{qr_adjacent_angles.png}

\paragraph{23. Перпендикуляр і похила.}
Спільна сторона ($OB$) двох суміжних кутів називається \textbf{похилою} до прямої ($AC$), на якій лежать дві інші сторони, у тому випадку, коли суміжні кути не рівні між собою (рис.~22, ліворуч); у тому ж випадку, коли суміжні кути рівні (рис.~22, праворуч) і коли, отже, кожен із кутів є прямим, спільна сторона називається \textbf{перпендикуляром} до прямої, на якій лежать дві інші сторони. Спільна вершина ($O$) у першому випадку називається \textbf{основою похилої}, у другому випадку — \textbf{основою перпендикуляра}.

\begin{center}
\begin{tikzpicture}[scale=0.95, >=Stealth]
    % Похила
    \draw[thick] (-2.5,0) node[below] {$A$} -- (2.0,0) node[below] {$C$};
    \draw[very thick, color=primaryblue, ->] (0,0) -- (1.5,2.0) node[above right] {$B$};
    \filldraw (0,0) circle (2pt) node[below] {$O$};
    \node at (0, -0.9) {\small\textbf{Похила $OB$}};

    % Перпендикуляр
    \begin{scope}[xshift=6cm]
        \draw[thick] (-2.0,0) node[below] {$A$} -- (2.0,0) node[below] {$C$};
        \draw[very thick, color=feynmangreen, ->] (0,0) -- (0,2.0) node[above] {$B$};
        \draw (0.3,0) -- (0.3,0.3) -- (0,0.3);
        \filldraw (0,0) circle (2pt) node[below] {$O$};
        \node at (0, -0.9) {\small\textbf{Перпендикуляр $OB$}};
    \end{scope}

    \node at (3.0, -1.6) {\small\textbf{Рис. 22.} Похила та перпендикуляр до прямої $AC$};
\end{tikzpicture}
\end{center}

Дві прямі ($AC$ і $BD$, рис.~21), що перетинаються між собою під прямим кутом, називаються \textbf{взаємно перпендикулярними}. Те, що пряма $AC$ перпендикулярна до прямої $BD$, виражають письмово так:
\[
AC \perp BD.
\]

\paragraph{24. Креслярський косинець.}
Для побудови перпендикуляра до даної прямої дуже зручним є креслярський косинець (трикутник), у якого один із кутів зроблено прямим. Щоб провести перпендикуляр до прямої $AB$ (рис.~23) через точку $C$, дану на цій прямій, або через точку $D$, взяту поза прямою, приставляють лінійку до прямої $AB$ і до лінійки косинець, і потім, притримуючи лінійку рукою, рухають косинець уздовж лінійки доти, доки друга сторона прямого кута не пройде через точку $C$ або $D$. Залишається потім провести пряму $CE$.

\begin{center}
\begin{tikzpicture}[scale=0.9, >=Stealth]
    % Лінійка
    \draw[fill=gray!20, thick] (-3,-0.5) rectangle (3.5,0);
    \draw[thick] (-3,0) -- (3.5,0);
    \node at (-2.7, 0.3) {$A$};
    \node at (3.2, 0.3) {$B$};

    % Косинець
    \draw[fill=primaryblue!15, very thick] (0,0) -- (0,3.0) -- (2.2,0) -- cycle;
    \draw (0,0.4) -- (0.4,0.4) -- (0.4,0);
    \filldraw (0,0) circle (2pt) node[below] {$C$};
    \filldraw (0,1.8) circle (2pt) node[left] {$D$};
    \node at (0, 3.3) {$E$};
\end{tikzpicture}\\[4pt]
{\small\textbf{Рис. 23.} Проведення перпендикуляра за допомогою косинця і лінійки}
\end{center}

\textit{Зауваження.}
\begin{enumerate}
    \item Якщо перпендикуляр до прямої $AB$ доводиться проводити з точки $C$, що лежить на цій прямій, то кажуть, що цей перпендикуляр треба \textit{«поставити»} (або провести) до прямої $AB$, а якщо потрібно перпендикуляр провести з точки $D$, що лежить поза прямою, то кажуть, що його треба \textit{«опустити»} на пряму (все одно: вниз, угору чи вбік).
    \item Очевидно, що через будь-яку точку даної прямої можна до цієї прямої провести перпендикуляр, і до того ж тільки один.
\end{enumerate}

\paragraph{25. Вертикальні кути.}
Два кути називаються \textbf{вертикальними}, якщо сторони одного становлять продовження сторін другого (є доповняльними променями). Так, при перетині двох прямих $AB$ і $CD$ (рис.~24) утворюються дві пари вертикальних кутів: $AOD$ і $COB$, $AOC$ і $DOB$ (і 4 пари суміжних кутів).

\begin{center}
\begin{tikzpicture}[scale=1.1, >=Stealth]
    \coordinate (O) at (0,0);
    \coordinate (A) at (-2.5,-1.2);
    \coordinate (B) at (2.5,1.2);
    \coordinate (C) at (-2,1.5);
    \coordinate (D) at (2,-1.5);

    \draw[thick, <->] (A) node[below left] {$A$} -- (B) node[above right] {$B$};
    \draw[thick, <->] (C) node[above left] {$C$} -- (D) node[below right] {$D$};
    \filldraw (O) circle (2pt) node[above=2pt] {$O$};

    \node at (-1.5, 0.2) {$\angle 1$};
    \node at (1.5, -0.2) {$\angle 2$};
    \node at (0, 1.0) {$\angle 3$};
    \node at (0, -1.0) {$\angle 4$};
\end{tikzpicture}\\[4pt]
{\small\textbf{Рис. 24.} Вертикальні кути: $\angle 1 = \angle 2$, $\angle 3 = \angle 4$}
\end{center}

\begin{theorem}[Властивість вертикальних кутів]
Два вертикальні кути рівні між собою:
\[
\angle AOD = \angle BOC \quad (\text{або } \angle 1 = \angle 2).
\]
\end{theorem}

\begin{proof}
Кожен із них є суміжним з одним і тим самим кутом (з $\angle DOB$ або з $\angle AOC$), а такі кути, як ми бачили у §~22, рівні один одному:
\begin{align*}
    \angle 1 + \angle 3 &= 180^\circ \implies \angle 1 = 180^\circ - \angle 3, \\
    \angle 2 + \angle 3 &= 180^\circ \implies \angle 2 = 180^\circ - \angle 3.
\end{align*}
Праві частини однакові, тому $\angle 1 = \angle 2$.
\end{proof}

\begin{feynmanbox}[Модель ножиць]
Уяви звичайні ножиці. Коли ти розкриваєш леза (кут 1), кільця ножиць з іншого боку гвинтика (кут 2) розкриваються рівно на той самий кут! 

Чому? Кожне лезо разом зі своєю ручкою утворює одну жорстку негнучку пряму лінію. Гвинтик — це точка перетину $O$. Якщо один кут збільшується, його суміжний партнер зменшується рівно на стільки ж, автоматично збільшуючи протилежний кут на ту саму величину.
\end{feynmanbox}

\motionlink{Вертикальні кути через півоберт}%
{Поверни копію кута навколо точки перетину на $180^\circ$. Його сторони сумістяться з протилежними променями, а сам кут --- із вертикальним.}%
{https://radchr.github.io/Algebra/geometry-vertical-half-turn.html}%
{qr_vertical_half_turn.png}

\paragraph{26. Кути, що мають спільну вершину.}
Про такі кути корисно помітити такі три прості істини:
\begin{enumerate}
    \item \textbf{Сума кількох кутів} ($AOB, BOC, COD, DOE$, рис.~25), що мають спільну вершину й заповнюють увесь простір по один бік від прямої ($AE$), \textbf{дорівнює $180^\circ$} (тобто двом прямим кутам, або $2d$), оскільки сума ця становить розгорнутий кут, а такий кут містить $180^\circ$.
    
    \begin{center}
    \begin{tikzpicture}[scale=0.9, >=Stealth]
        \coordinate (O) at (0,0);
        \draw[thick] (-3,0) node[left] {$A$} -- (3,0) node[right] {$E$};
        \draw[thick, ->] (O) -- (-2,1.8) node[above left] {$B$};
        \draw[thick, ->] (O) -- (0,2.2) node[above] {$C$};
        \draw[thick, ->] (O) -- (2,1.8) node[above right] {$D$};
        \filldraw (O) circle (2pt) node[below] {$O$};
    \end{tikzpicture}\\[4pt]
    {\small\textbf{Рис. 25.} Кути по один бік від прямої: сума $= 180^\circ$}
    \end{center}

    \item \textbf{Сума кількох кутів} ($AOB, BOC, COD, DOE, EOA$, рис.~26), що мають спільну вершину й заповнюють увесь простір навколо цієї вершини, \textbf{дорівнює $360^\circ$} (тобто чотирьом прямим кутам, або $4d$), оскільки ця сума утворює повний кут, а такий кут містить $360^\circ$.
    
    \begin{center}
    \begin{tikzpicture}[scale=0.85, >=Stealth]
        \coordinate (O) at (0,0);
        \draw[thick, ->] (O) -- (2.2,0) node[right] {$A$};
        \draw[thick, ->] (O) -- (1.2,2.0) node[above right] {$B$};
        \draw[thick, ->] (O) -- (-1.5,1.8) node[above left] {$C$};
        \draw[thick, ->] (O) -- (-2.2,-0.5) node[left] {$D$};
        \draw[thick, ->] (O) -- (0.3,-2.0) node[below] {$E$};
        \filldraw (O) circle (2pt) node[below left=2pt] {$O$};
    \end{tikzpicture}\\[4pt]
    {\small\textbf{Рис. 26.} Кути навколо спільної вершини: сума $= 360^\circ$}
    \end{center}

    \item \textbf{Якщо два кути} ($AOB$ і $BOC$, рис.~27) мають спільну вершину ($O$) і спільну сторону ($OB$) і в сумі становлять $180^\circ$ (два прямих кути, або $2d$), то їхні дві інші сторони ($OA$ і $OC$) становлять продовження одна одної (тобто такі кути будуть суміжними). Справді, сума вказаних кутів становить кут $AOC$, і якщо цей кут дорівнює $180^\circ$, то це означає, що він є розгорнутим кутом, тобто що $AO$ становить із $OC$ одну пряму.
    
    \begin{center}
    \begin{tikzpicture}[scale=1.0, >=Stealth]
        \coordinate (O) at (0,0);
        \coordinate (A) at (-3,0);
        \coordinate (C) at (3,0);
        \coordinate (B) at (1.2,2.2);

        \draw[thick] (A) node[below] {$A$} -- (C) node[below] {$C$};
        \draw[very thick, color=primaryblue, ->] (O) -- (B) node[above right] {$B$};
        \filldraw (O) circle (2pt) node[below] {$O$};
    \end{tikzpicture}\\[4pt]
    {\small\textbf{Рис. 27.} Ознака суміжних кутів}
    \end{center}
\end{enumerate}

\clearpage
\section{Вправи}

\begin{problem}
Деякий кут дорівнює $38^\circ 29'$. Знайти величину суміжного з ним кута.
\end{problem}

\begin{center}
\begin{tikzpicture}[scale=0.95, >=Stealth]
    \coordinate (O) at (0,0);
    \coordinate (A) at (-3,0);
    \coordinate (C) at (3,0);
    \coordinate (B) at (38.5:2.6);
    
    \draw[thick] (A) node[below] {$A$} -- (C) node[below] {$C$};
    \draw[very thick, color=primaryblue, ->] (O) -- (B) node[above right] {$B$};
    \filldraw (O) circle (2pt) node[below] {$O$};
    
    % Кут 38°29'
    \draw[thick, color=primaryblue] (0.8,0) arc (0:38.5:0.8);
    \node[color=primaryblue, right] at (0.9,0.3) {\small $38^\circ 29'$};
    
    % Шуканий кут ?
    \draw[thick, color=warningred] (-0.9,0) arc (180:38.5:0.9);
    \node[color=warningred, left] at (-0.7,0.7) {\small\textbf{?}};
\end{tikzpicture}\\[4pt]
{\small\textbf{Рис. до вправи 1.} Знаходження суміжного кута}
\end{center}

\textbf{Розв'язання:}
Оскільки сума двох суміжних кутів дорівнює $180^\circ$, то шуканий кут дорівнює:
\[
180^\circ - 38^\circ 29' = 179^\circ 60' - 38^\circ 29' = 141^\circ 31'.
\]
\textbf{Відповідь:} $141^\circ 31'$.

\vspace{0.4cm}

\begin{problem}
Два кути $ABC$ і $CBD$, маючи спільну вершину $B$ і спільну сторону $BC$, розташовані так, що вони не покривають один одного: кут $ABC = 100^\circ 20'$, а кут $CBD = 79^\circ 40'$. Чи становлять сторони $AB$ і $BD$ пряму чи ламану?
\end{problem}

\begin{center}
\begin{tikzpicture}[scale=0.95, >=Stealth]
    \coordinate (B) at (0,0);
    \coordinate (D) at (2.8,0);
    \coordinate (A) at (-2.8,0);
    \coordinate (C) at (79.67:2.4);
    
    \draw[dashed, color=gray] (A) -- (D);
    \draw[thick, ->] (B) -- (A) node[below] {$A$};
    \draw[thick, ->] (B) -- (D) node[below] {$D$};
    \draw[very thick, color=primaryblue, ->] (B) -- (C) node[above] {$C$ (спільна сторона)};
    \filldraw (B) circle (2pt) node[below] {$B$};
    
    \draw[thick, color=feynmangreen] (0.8,0) arc (0:79.67:0.8);
    \node[color=feynmangreen] at (40:1.3) {\small $79^\circ 40'$};
    
    \draw[thick, color=primaryblue] (79.67:0.7) arc (79.67:180:0.7);
    \node[color=primaryblue] at (130:1.2) {\small $100^\circ 20'$};
\end{tikzpicture}\\[4pt]
{\small\textbf{Рис. до вправи 2.} Сторони $AB$ і $BD$ утворюють розгорнутий кут ($180^\circ$)}
\end{center}

\textbf{Розв'язання:}
Знайдемо суму цих кутів:
\[
\angle ABC + \angle CBD = 100^\circ 20' + 79^\circ 40' = 179^\circ 60' = 180^\circ.
\]
Оскільки кути мають спільну вершину $B$, спільну сторону $BC$, не перекриваються і їхня сума дорівнює $180^\circ$, то за властивістю §~26 (пункт 3) кут $\angle ABD$ є розгорнутим, тобто промені $BA$ і $BD$ є доповняльними. Отже, сторони $AB$ і $BD$ становлять \textbf{пряму лінію}.

\vspace{0.4cm}

\begin{problem}
Побудувати будь-який кут і за допомогою транспортира та лінійки провести його бісектрису.
\end{problem}

\begin{center}
\begin{tikzpicture}[scale=0.95, >=Stealth]
    \coordinate (O) at (0,0);
    \coordinate (A) at (3,0);
    \coordinate (B) at (60:2.8);
    \coordinate (K) at (30:3.1);
    
    \draw[thick, ->] (O) -- (A) node[right] {$A$};
    \draw[thick, ->] (O) -- (B) node[above right] {$B$};
    \draw[very thick, color=feynmangreen, ->] (O) -- (K) node[right] {$K$ (бісектриса)};
    \filldraw (O) circle (2pt) node[left] {$O$};
    
    \draw[thick, color=feynmangreen] (1.0,0) arc (0:30:1.0);
    \node[color=feynmangreen] at (15:1.4) {\small $30^\circ$};
    
    \draw[thick, color=feynmangreen] (30:1.2) arc (30:60:1.2);
    \node[color=feynmangreen] at (45:1.6) {\small $30^\circ$};
\end{tikzpicture}\\[4pt]
{\small\textbf{Рис. до вправи 3.} Поділ кута $60^\circ$ бісектрисою на два кути по $30^\circ$}
\end{center}

\textbf{Вказівка:} Виміряйте кут транспортиром, поділіть отримане число градусів на 2, позначте точку на шкалі транспортира, що відповідає половині кута, і через цю точку та вершину проведіть промінь лінійкою.

\vspace{0.4cm}

\begin{problem}[Класична теорема про бісектриси суміжних кутів]
Довести, що бісектриси двох суміжних кутів взаємно перпендикулярні.
\end{problem}

\begin{center}
\begin{tikzpicture}[scale=1.05, >=Stealth]
    \coordinate (O) at (0,0);
    \coordinate (A) at (-3,0);
    \coordinate (C) at (3,0);
    \coordinate (B) at (60:2.8);
    
    % Бісектриси: 30° та 120°
    \coordinate (K) at (30:2.9);
    \coordinate (L) at (120:2.9);
    
    \draw[thick] (A) node[below] {$A$} -- (C) node[below] {$C$};
    \draw[thick, dashed] (O) -- (B) node[above right] {$B$ (спільна сторона)};
    \draw[very thick, color=primaryblue, ->] (O) -- (K) node[above right] {$OK$ (бісектриса)};
    \draw[very thick, color=feynmangreen, ->] (O) -- (L) node[above left] {$OL$ (бісектриса)};
    \filldraw (O) circle (2pt) node[below] {$O$};
    
    % Прямий кут між бісектрисами (30° і 120° -> різниця 90°)
    \draw[thick, color=warningred] (30:0.6) -- ++(120:0.6) -- ++(210:0.6);
    \node[color=warningred, font=\bfseries] at (75:0.95) {$90^\circ$};
    
    % Позначки рівних кутів
    \draw[color=primaryblue] (0.9,0) arc (0:30:0.9) node[midway, right=1pt] {\tiny $\frac{\alpha}{2}$};
    \draw[color=primaryblue] (30:1.1) arc (30:60:1.1) node[midway, right=1pt] {\tiny $\frac{\alpha}{2}$};
    
    \draw[color=feynmangreen] (60:1.2) arc (60:120:1.2) node[midway, above] {\tiny $\frac{\beta}{2}$};
    \draw[color=feynmangreen] (120:1.0) arc (120:180:1.0) node[midway, left=1pt] {\tiny $\frac{\beta}{2}$};
\end{tikzpicture}\\[4pt]
{\small\textbf{Рис. до вправи 4.} Бісектриси $OK \perp OL$, оскільки $\frac{\alpha}{2} + \frac{\beta}{2} = 90^\circ$}
\end{center}

\textbf{Доведення:}
Нехай дані суміжні кути мають градусні міри $\alpha$ та $\beta$.
За теоремою про суму суміжних кутів:
\[
\alpha + \beta = 180^\circ.
\]
Бісектриса першого кута утворює з його спільною стороною кут $\frac{\alpha}{2}$. Бісектриса другого кута утворює з тією самою спільною стороною кут $\frac{\beta}{2}$.
Кут між двома бісектрисами дорівнює сумі цих половин:
\[
\text{Кут між бісектрисами} = \frac{\alpha}{2} + \frac{\beta}{2} = \frac{\alpha + \beta}{2} = \frac{180^\circ}{2} = 90^\circ.
\]
Кут, рівний $90^\circ$, є прямим. Отже, бісектриси взаємно перпендикулярні. Що й треба було довести.

\motionlink{Бісектриси суміжних кутів}%
{Змінюй суміжні кути й перевіряй, що кут між їхніми бісектрисами незмінно дорівнює $90^\circ$.}%
{https://radchr.github.io/Algebra/geometry-adjacent-bisectors.html}%
{qr_adjacent_bisectors.png}

\vspace{0.4cm}

\begin{problem}
Довести, що бісектриси двох вертикальних кутів становлять продовження одна одної (лежать на одній прямій).
\end{problem}

\begin{center}
\begin{tikzpicture}[scale=1.0, >=Stealth]
    \coordinate (O) at (0,0);
    \coordinate (A) at (-2.6,-1.3);
    \coordinate (B) at (2.6,1.3);
    \coordinate (C) at (-2.6,1.3);
    \coordinate (D) at (2.6,-1.3);
    
    \draw[thick] (A) node[below left] {$A$} -- (B) node[above right] {$B$};
    \draw[thick] (C) node[above left] {$C$} -- (D) node[below right] {$D$};
    \filldraw (O) circle (2pt) node[below=3pt] {$O$};
    
    % Бісектриси вертикальних кутів
    \draw[very thick, color=warningred, dashed, <->] (-3.1,0) node[left] {$L$} -- (3.1,0) node[right] {$K$};
    
    % Дужки кутів
    \draw[color=primaryblue] (0.8,0) arc (0:26.5:0.8);
    \draw[color=primaryblue] (0.8,0) arc (0:-26.5:0.8);
    \node at (1.4, 0.25) {\tiny $\frac{\alpha}{2}$};
    \node at (1.4, -0.25) {\tiny $\frac{\alpha}{2}$};
    
    \draw[color=primaryblue] (-0.8,0) arc (180:153.5:0.8);
    \draw[color=primaryblue] (-0.8,0) arc (180:206.5:0.8);
    \node at (-1.4, 0.25) {\tiny $\frac{\alpha}{2}$};
    \node at (-1.4, -0.25) {\tiny $\frac{\alpha}{2}$};
\end{tikzpicture}\\[4pt]
{\small\textbf{Рис. до вправи 5.} Бісектриси $OK$ та $OL$ утворюють розгорнутий кут $180^\circ$ (пряму лінію)}
\end{center}

\textbf{Доведення:}
Нехай при перетині двох прямих утворилися вертикальні кути $\angle 1$ і $\angle 2$, причому $\angle 1 = \angle 2 = \alpha$, а кут $\angle 3$ є суміжним з кожним із них ($\angle 3 = 180^\circ - \alpha$).
Проведемо бісектриси вертикальних кутів. Бісектриса кута $\angle 1$ ділить його на два кути по $\frac{\alpha}{2}$. Бісектриса кута $\angle 2$ також ділить його на два кути по $\frac{\alpha}{2}$.
Кут між цими двома бісектрисами складається з половини кута $\angle 1$, суміжного кута $\angle 3$ і половини кута $\angle 2$:
\[
\text{Кут} = \frac{\alpha}{2} + \angle 3 + \frac{\alpha}{2} = \alpha + (180^\circ - \alpha) = 180^\circ.
\]
Оскільки кут між бісектрисами дорівнює $180^\circ$, він є розгорнутим, отже, його сторони (бісектриси) лежать на одній прямій і є доповняльними променями. Що й треба було довести.

\motionlink{Бісектриси вертикальних кутів}%
{Рухай одну з прямих і спостерігай: бісектриси протилежних кутів завжди залишаються протилежними променями.}%
{https://radchr.github.io/Algebra/geometry-vertical-bisectors.html}%
{qr_vertical_bisectors.png}

\vspace{0.4cm}

\begin{problem}
Якщо при точці $O$ прямої $AB$ (рис.~24) побудуємо по різні боки від $AB$ рівні кути $AOD$ і $BOC$, то сторони їх $OD$ і $OC$ становлять одну пряму.
\end{problem}

\begin{center}
\begin{tikzpicture}[scale=0.95, >=Stealth]
    \coordinate (O) at (0,0);
    \coordinate (A) at (-3,0);
    \coordinate (B) at (3,0);
    \coordinate (C) at (35:2.6);
    \coordinate (D) at (215:2.6);
    
    \draw[thick] (A) node[below] {$A$} -- (B) node[below] {$B$};
    \draw[very thick, color=primaryblue, ->] (O) -- (C) node[above right] {$C$};
    \draw[very thick, color=primaryblue, ->] (O) -- (D) node[below left] {$D$};
    \draw[dashed, color=gray] (D) -- (C);
    \filldraw (O) circle (2pt) node[below right=2pt] {$O$};
    
    \draw[color=primaryblue, thick] (0.8,0) arc (0:35:0.8);
    \node[color=primaryblue] at (17.5:1.2) {\small $\alpha$};
    
    \draw[color=primaryblue, thick] (-0.8,0) arc (180:215:0.8);
    \node[color=primaryblue] at (197.5:1.2) {\small $\alpha$};
\end{tikzpicture}\\[4pt]
{\small\textbf{Рис. до вправи 6.} $\angle AOD = \angle BOC = \alpha \implies COD$ — пряма лінія}
\end{center}

\textbf{Доведення:}
Кути $BOC$ і $AOC$ є суміжними, оскільки точки $A, O, B$ лежать на одній прямій, тому:
\[
\angle AOC + \angle BOC = 180^\circ.
\]
За умовою $\angle BOC = \angle AOD$. Замінимо в рівності кут $\angle BOC$ на рівний йому $\angle AOD$:
\[
\angle AOC + \angle AOD = 180^\circ.
\]
Оскільки кути $AOC$ і $AOD$ мають спільну вершину $O$, спільну сторону $OA$, лежать по різні боки від прямої $CD$ і в сумі дають $180^\circ$, то за ознакою §~26 (пункт 3) промені $OC$ і $OD$ становлять одну пряму лінію. Що й треба було довести.

\vspace{0.4cm}

\begin{problem}
Якщо з точки $O$ (рис.~24) проведемо промені $OA, OD, OB, OC$ так, що $\angle AOC = \angle DOB$ і $\angle AOD = \angle COB$, то $OB$ є продовження $OA$, а $OD$ — продовження $OC$.
\end{problem}

\begin{center}
\begin{tikzpicture}[scale=0.95, >=Stealth]
    \coordinate (O) at (0,0);
    \coordinate (A) at (-2.8,0);
    \coordinate (B) at (2.8,0);
    \coordinate (C) at (50:2.4);
    \coordinate (D) at (230:2.4);
    
    \draw[thick, ->] (O) -- (A) node[left] {$A$};
    \draw[thick, ->] (O) -- (B) node[right] {$B$};
    \draw[thick, ->] (O) -- (C) node[above right] {$C$};
    \draw[thick, ->] (O) -- (D) node[below left] {$D$};
    \filldraw (O) circle (2pt) node[below=3pt] {$O$};
    
    % Дужки кутів 1 і 2
    \draw[color=primaryblue, thick] (0.7,0) arc (0:50:0.7) node[midway, right=1pt] {\small $\alpha$};
    \draw[color=primaryblue, thick] (-0.7,0) arc (180:230:0.7) node[midway, left=1pt] {\small $\alpha$};
    
    % Дужки кутів 3 і 4
    \draw[color=warningred, thick] (50:0.9) arc (50:180:0.9) node[midway, above] {\small $\beta$};
    \draw[color=warningred, thick] (230:0.9) arc (230:360:0.9) node[midway, below] {\small $\beta$};
\end{tikzpicture}\\[4pt]
{\small\textbf{Рис. до вправи 7.} Дві пари рівних протилежних кутів: $\alpha + \beta = 180^\circ$}
\end{center}

\textbf{Доведення:}
Сума всіх чотирьох кутів навколо спільної вершини $O$ за §~26 (пункт 2) дорівнює $360^\circ$:
\[
\angle AOC + \angle COB + \angle DOB + \angle AOD = 360^\circ.
\]
Враховуючи, що $\angle AOC = \angle DOB$ та $\angle COB = \angle AOD$, отримуємо:
\[
2 \angle AOC + 2 \angle COB = 360^\circ \implies \angle AOC + \angle COB = 180^\circ.
\]
Оскільки кути $AOC$ і $COB$ мають спільну сторону $OC$, а їхня сума дорівнює $180^\circ$, то промені $OA$ і $OB$ утворюють розгорнутий кут, тобто $OB$ є продовженням $OA$.
Аналогічно:
\[
\angle AOC + \angle AOD = 180^\circ \implies OD \text{ є продовженням } OC.
\]
Що й треба було довести.
